\section{Conclusion}
\label{sec:conclusion}

We studied whether prompt-only frontier LLMs can recommend music from listening history or playlist context without any task-specific training. Using a shared closed-set ranking protocol on \lastfm and \mpd, we found that both \gptfourone and \claudesonnet outperform lightweight classical baselines by large margins. The best system, \gptfourone with a minimal context prompt, reaches 0.589 \mrr on \lastfm and 0.329 \mrr on \mpd, well above the strongest implemented baseline.

The key takeaway is simple: prompt-only LLM recommendation can work surprisingly well against lightweight public baselines, especially when the task is framed as reranking a fixed candidate pool from recent context. The richer profile prompt does not help, and the two frontier models are statistically tied in the strongest setting.

The main next steps are clear. Future work should compare against stronger public baselines such as SASRec, PureSVD, or RecBole models; separate candidate generation from LLM reranking; evaluate real personal Spotify exports; and add human judgments of novelty, usefulness, and playlist fit. Those steps are needed before making any stronger claim about performance relative to production music recommenders.
